\chapter*{Remerciements}
\addcontentsline{toc}{chapter}{Remerciements}

La réalisation de ce projet de fin d'études n'aurait pas été possible sans le concours
bienveillant de nombreuses personnes à qui nous tenons à exprimer notre sincère gratitude.

Nous adressons en premier lieu nos remerciements les plus chaleureux à notre encadrante
académique, \textbf{Mme Najiba TAGOUGUI}, pour la qualité de son encadrement, la
pertinence de ses conseils scientifiques et sa disponibilité constante tout au long de
ce travail. Sa rigueur intellectuelle et son soutien méthodologique ont été déterminants
dans la conduite et la réussite de ce projet.

Nous remercions également notre encadrant professionnel, \textbf{M. HoussemEddine
HASSAYOUNE}, ingénieur chez Mobelite, pour nous avoir accueillis au sein de son équipe,
partagé son expertise terrain en matière de sécurité applicative et orienté nos choix
techniques avec une vision résolument pragmatique. Sa connaissance des enjeux DevSecOps
a été une source d'inspiration continue.

Nos remerciements vont aussi à la direction et aux équipes de \textbf{Mobelite} pour
avoir mis à notre disposition un environnement de travail stimulant, un accès aux
ressources nécessaires et un cadre professionnel propice à l'innovation.

Nous tenons à remercier les membres du jury, qui ont accepté d'évaluer ce travail et
dont les remarques contribueront à son enrichissement. Leur expertise représente pour
nous une reconnaissance précieuse de l'effort fourni.

Nous exprimons notre profonde gratitude à l'ensemble du corps enseignant de
l'\textbf{Institut Supérieur d'Informatique et de Multimédia de Sfax}, dont la
formation solide nous a dotés des compétences techniques et analytiques nécessaires
pour mener à bien un projet de cette envergure.

Enfin, nous remercions nos familles respectives pour leur soutien moral indéfectible,
leur patience et leur encouragement tout au long de ces années d'études. Ce travail
leur appartient autant qu'à nous.

\clearpage